\documentclass[12pt]{article}

%% Basic document formatting
\usepackage{amsmath, amsthm, amssymb, amsfonts}
\usepackage{mathrsfs}
\usepackage{mathtools}
\usepackage{xspace}
\usepackage{thmtools}
\usepackage{graphicx}
\usepackage{tikz}
\usepackage{tikz-cd} 
\usepackage{ytableau}
\usepackage{blkarray}
\usepackage{hyperref}
\usepackage{caption}
\usepackage{setspace}
\usepackage{array}
\usepackage{fancyhdr}
\usepackage{titling}
\usepackage[left=0.4in,right=0.4in,top=1in,bottom=1in]{geometry}
\usepackage{float}
\usepackage{cancel}
\usepackage{tabularx}
\usepackage[utf8]{inputenc}
\usepackage[english]{babel}
\usepackage{framed}
\usepackage[dvipsnames]{xcolor}
\usepackage{environ}
\usepackage{tcolorbox}
\tcbuselibrary{theorems,skins,breakable}

\usepackage{enumitem}
\setlist[enumerate]{leftmargin=*}
\setlist[enumerate,1]{labelindent=\parindent}
\setlist[enumerate,2]{labelindent=0pt}

% Blackboard Bold
\newcommand{\Z}{\mathbb{Z}}                 % integers
\newcommand{\Zpl}{\mathbb{Z}_{+}}           % positive integers
\newcommand{\Znn}{\mathbb{Z}_{\geq 0}}      % nonnegative integers. 
\newcommand{\Q}{\mathbb{Q}}                 % rationals
\newcommand{\Qpl}{\mathbb{Q}_{+}}           % positive Rationals
\newcommand{\R}{\mathbb{R}}                 % reals
\newcommand{\Rpl}{\mathbb{R}_{+}}           % positive Reals
\newcommand{\C}{\mathbb{C}}                 % complex numbers
\newcommand{\F}{\mathbb{F}}                 % field

% Words
\newcommand{\st}{\text{ such that }}
\newcommand{\wrt}{\text{ with respect to }}
\newcommand{\with}{\text{ with }}
\newcommand{\ie}{\text{, i.e. }}
\newcommand*{\ditto}{\text{---}\texttt{"}\text{---}}

% Operators
\newcommand{\abs}[1]{\left|#1\right|}                   % absolute value
\newcommand{\norm}[1]{\abs{\abs{#1}}}
\newcommand{\floor}[1]{\left\lfloor #1 \right\rfloor}   % floor
\newcommand{\ceil}[1]{\left\lceil #1 \right\rceil}      % ceiling
\newcommand{\powset}[1]{\mathscr{P}(#1)}

% Category Theory 
\renewcommand{\hom}{\text{Hom}}
\newcommand{\homspace}[3]{\hom_{#1}(#2, #3)}
\newcommand{\coproduct}[9]{
  \begin{tikzcd}[ampersand replacement=\&]
      \& \& #1 \arrow[dll, bend right, "#6"] \arrow[dl, "#5"] \\
   #4 \& #3 \arrow[l, "#9"] \\
      \& \& #2 \arrow[ull, bend left, "#8"] \arrow[ul, "#7"]
  \end{tikzcd}
}

\newcommand{\product}[9]{ 
  \begin{tikzcd}[ampersand replacement=\&]
      \& \& #3 \\
      #1 \arrow[r, "#7"] \arrow[urr, bend left, "#5"] \arrow[drr, bend right, "#6"] \& #2 \arrow[ur, "#8"] \arrow[dr, "#9"] \\ 
      \& \& #4
  \end{tikzcd}
}


% Algebra 
\newcommand{\Syl}[2]{\text{Syl}_{#1}{(#2)}}           % set of Sylow-p subgroups 
\newcommand{\<}{\langle}                % \<x\>, subgroup generated by x 
\renewcommand{\>}{\rangle}
\renewcommand{\index}[2]{[#1 : #2]}
\newcommand{\id}{\text{id}}             % identity element
\newcommand{\lhdneq}{%
  \mathrel{\ooalign{$\lneq$\cr\raise.22ex\hbox{$\lhd$}\cr}}}
\newcommand{\order}[1]{\text{o}(#1)}    % order of an element
\newcommand{\spec}{\operatorname{Spec}}
\newcommand{\rad}{\operatorname{rad}}   % radical of an ideal 
\newcommand{\sym}[1]{\mathfrak{S}_{#1}}
\newcommand{\aut}[1]{\text{Aut}(#1)} 
\newcommand{\gal}[2]{\text{Gal}(#1 / #2)} 
\newcommand{\inn}[1]{\text{Inn}(#1)} 
\let\oldcong\cong
\let\oldequiv\equiv
\renewcommand{\cong}{\oldequiv}
\renewcommand{\equiv}{\oldcong}

\newcommand{\triangleleftneq}{\mathrel{\ooalign{%
  $\trianglelefteq$\cr
  \hidewidth\raise0.06ex\hbox{$\not\mkern4mu$}\hidewidth\cr
}}}

\makeatletter % cycle
\newcommand{\cyc}[1]{(\mathbf{\cyc@process#1\relax})}
\def\cyc@process#1#2\relax{%
	#1%
	\ifx\relax#2\relax
	\else
		\,\cyc@process#2\relax
	\fi
}
\makeatother

\newcommand{\GL}[2]{\text{GL}_{#1}(#2)}                             % general linear group
\newcommand{\SL}[2]{\text{SL}_{#1}(#2)}                             % special linear group
\newcommand{\gl}[2]{\mathfrak{gl}_{#1}(#2)}
\renewcommand{\sl}[2]{\mathfrak{sl}_{#1}(#2)}
\newcommand{\so}[2]{\mathfrak{so}_{#1}(#2)}
\newcommand{\su}[1]{\mathfrak{su}(#1)}

% Linear Algebra
\newcommand{\bmat}[1]{\begin{bmatrix}#1\end{bmatrix}}   % bracketed matrix
\newcommand{\rank}{\operatorname{rank}}                 % rank
\newcommand{\nullity}{\operatorname{nullity}}           % nullity
\newcommand{\tr}{\operatorname{tr}}

% Combinatorics
\newcommand{\qnum}[1]{\left[#1\right]_q}                % q-analog
\newcommand{\qfact}[1]{\left[#1\right]!_q}              % q-analog factorial 
\newcommand{\qbinom}[2]{\binom{#1}{#2}_q}               % Gaussian binomial coefficient

% Topology/Analysis
\newcommand{\ball}[2]{\text{B}_{#1}(#2)}  % B_r(x): open r-balls around x
\newcommand{\diam}{\text{diam}}           % diamter of a set in metric space 
\newcommand{\lowint}[2]{\underline{\int_{#1}^{#2}}}
\newcommand{\upint}[2]{\overline{\int}_{#1}^{#2}}

% Misc Notation
\newcommand{\oneton}[1]{\{1, \ldots, #1\}}
\newcommand{\defeq}{\vcentcolon=}   % :=
\newcommand{\eqdef}{=\vcentcolon}   % =: 
\renewcommand{\bf}[1]{\textbf{#1}}
\newcommand{\im}{\operatorname{im}}
\newcommand{\fnrest}[2]{\left.#1\right|_{#2}}
\newcommand{\isom}{\overset{\sim}{\rightarrow}} 


\renewcommand{\thesection}{\S \arabic{section}}
\renewcommand{\thesubsection}{\arabic{subsection}.}
\renewcommand{\thesubsubsection}{(\alph{subsubsection})} 
\newtheorem{claim}{Claim}
\newtheorem*{lemma}{Lemma}

%% Headers & title setup
\newcommand{\course}{\textit{Algebra: Chapter 0}, Aluffi} 
\newcommand{\myname}{Jay Ser}
\setlength{\headheight}{14.5pt}
\pagestyle{fancy}
\fancyhf{}
\renewcommand{\headrulewidth}{0.4pt}
\lhead{\course}
\rhead{\myname}
\cfoot{\thepage}
\setlength{\droptitle}{-4em} 
\title{\course\ \\ Part 6: Linear Algebra}
\author{\myname}
\date{May 25, 2026 $\sim$ \today}

\begin{document}
\maketitle
\thispagestyle{fancy}

\newcommand{\dirsum}[2]{#1^{\oplus#2}}
\newcommand{\rmod}[1]{#1\text{-module}\xspace}

%------------------------------------------------------------------------------%

\section{Free Modules Revisited}
\setcounter{subsection}{1}
\subsection{} 
$\tr(aM) = a\tr(M)$, so $\sl{n}{\R}$ and $\sl{n}{\C}$ are $\R$-vector spaces by the natural action $a \cdot M \defeq aM$. 
The following gives an $\R$-basis for $\sl{n}{\R}$: 
\begin{equation} \label{eq:1_2_slnRbasis}
  \{e_{kl} \mid k \neq l\} \cup \{e_{11} - e_{kk} \mid k > 1\}
\end{equation}
where $e_{kl}$ is the $n \times n$ matrix with a 1 in the $(k, l)$th entry and zeros everywhere else. 
It is clear that the elements of \eqref{eq:1_2_slnRbasis} are linearly independent. 
It is also clear $\{e_{kl} \mid k \neq l\}$ span the nondiagonal entries of all matrices in $\sl{n}{\R}$. 
The following shows that 
\begin{equation} \label{eq:1_2_diagspan}
  \{e_{11} - e_{kk} \mid k > 1\}
\end{equation}
span the diagonal entries of all the elements in $\sl{n}{\R}$, which completes the verification that \eqref{eq:1_2_slnRbasis} spans $\sl{n}{\R}$. 
Induct on $n$: 
for $n = 2$, it is clear that $\bmat{1 \\ & -1}$ spans all possible diagonal entries for the elements of $\sl{2}{\R}$. 
Suppose the proposition is true for $n - 1$. 
Take any $M = x_1 e_{11} + x_2 e_{22} + \ldots + x_n e_{nn} \in \sl{n}{\R}$; 
that is $x_1 + x_2 + \ldots + x_n = 0$.
Then the manipulation
\begin{equation} \label{eq:1_2_Mmanip}
M = \bmat{x_1 + x_n \\ & x_2 \\ & & \ddots \\ & & & x_{n - 1} \\ & & & & 0} - x_n \bmat{1 \\ &  0 \\ & & \ddots \\ & & & 0 \\ & & & & -1},
\end{equation}
the natural injection of $\sl{n - 1}{\R}$ into the first $(n - 1)$ rows and columns of $\sl{n}{R}$, and the observation that the $(n - 1) \times (n - 1)$ submatrix of the first matrix on the right hand side of \eqref{eq:1_2_Mmanip} is in $\sl{n - 1}{\R}$ shows that $M$ is generated by \eqref{eq:1_2_diagspan}.

Hence \eqref{eq:1_2_slnRbasis} is a basis for $\sl{n}{\R}$; 
accordingly, 
$$\dim{\sl{n}{\R}} = n^2 - n + (n - 1) = n^2 - 1.$$
The basis for $\sl{n}{\C}$ is given similarly, except one has to remember $\C$ is a two-dimensional vector space over $\R$: 
\begin{gather*}
  \{e_{kl} \mid k \neq l\} \cup \{i e_{kl} \mid k \neq l\} \cup \{e_{11} - e_{kk} \mid k > 1\} \cup \{ie_{11} - ie_{kk} \mid k > 1\} \text{ gives a basis,}\\ 
  \dim{\sl{n}{\C}} = 2n^2 - 2
\end{gather*}
where $i \defeq \sqrt{-1}$. 

Next, $a(M + M^{t}) = aM + aM^{t}$ shows that $\so{n}{\R}$ is also $\R$-vecctor spaces by the natural action $a \cdot M \defeq aM$. 
$M + M^{t} = 0 \iff a_{kl} = -a_{lk}$ for all $k, l$; 
namely, $M + M^{t} = 0 \implies a_{kk} = 0$ for all $k$, so any real matrix that satisfies $M + M^{t} = 0$ is a member of $\sl{n}{\R}$. 
$$\{e_{kl} - e_{lk} \mid k < l\}$$ 
gives a basis for $\so{n}{\R}$. 
Accordingly, 
$$\dim{\so{n}{\R}} = \sum_{k = 1}^{n} \sum_{l = k + 1}^{n} 1 = \frac{n(n - 1)}{2}$$

Proceed similarly with $\su{n}$. 
The natural action $a \cdot M \defeq aM$ makes $\su{n}$ an $\R$-vector space. 
$M + M^{*} = 0 \iff a_{kl} = -\overline{a_{lk}}$. 
Particularly, each $a_{kk}$ only has a complex part, and one mustn't forget that $\su{n}$ is a subset of $\sl{n}{\C}$.  
Hence
$$\{e_{kl} - e_{lk} \mid k < l\} \cup \{ie_{kl} + ie_{lk} \mid k < l\} \cup \{ie_{11} - ie_{kk} \mid k > 1\}$$ 
gives a basis for $\su{n}$ over $\R$, and 
$$\dim{\su{n}} = n(n - 1) + (n - 1) = (n + 1)(n - 1)$$

\setcounter{subsection}{4} 

\subsection{} 
\subsubsection{} 
This is false. 
As mentioned in the text, $\{2\}$ doesn't is a maximal linearly independent subset of $\Z$ as a free module over itself, but it doesn't generate $\Z$. 
Hence $\{2\}$ can never be completed into a basis for $\Z$; 
adding any other element will make the set linearly dependent.  

\subsubsection{} 
By a similar argument as part (a), this is false. 
$\{2, 3\}$ generate $\Z$ as a module over itself because the two numbers are coprime. 
However, as above, neither $\{2\}$ nor $\{3\}$ are a basis for $\Z$; 
namely, they don't generate all of $\Z$. 

\setcounter{subsection}{6}
\subsection{} 
Assume $S \subset M$ is linearly independent over $R$.
Suppose $S$ is not linearly independent as a subset of $V$ over $K$: 
there are $c_\alpha \in K$, finitely many of which aren't zero, and 
$$\sum c_\alpha s_\alpha = 0$$
Since only finitely many $c_\alpha$ aren't zero, one can cancel denominators: 
there is some $a \in R$ such that $a c_\alpha \in R \setminus 0$ for all $\alpha$ such that $c_\alpha \neq 0$: 
$$\sum ac_\alpha s_\alpha = 0, \quad (ac_\alpha \in R \setminus 0 \text{ for } \alpha \st c_\alpha \neq 0)$$
This gives a nontrivial $R$-linear dependence relation of elements of $S$, contradicting the assumption that $S \subset M$ is linearly independent over $R$. 
Hence $S$ must be a linearly independent subset of $V$ over $K$. 

The reverse direction is clear. 
Hence a maximal linearly independent subset of $M$ is a maximal linearly independent subset of $V$. 
A maximal linearly independent subset of $V$ is a basis for $V$, so if $B$ is such a set, 
$$\rank{M} = \#B = \dim{V}$$

Suppose $S$ generates $M$ over $R$. 
That is, for all $\sum r_\alpha a_\alpha \in M$,
$$\sum r_\alpha a_\alpha = \sum r'_\beta s_\beta$$ 
for some $r'_\beta \in R$ with at most finitely many of them nonzero. 
Given any $\sum c_\alpha a_\alpha \in V$, one can cancel the denominators of the nonzero coefficients to get $\sum c_\alpha a_\alpha = \frac{1}{b} \sum bc_\alpha a_\alpha$ with $b \in R$ and $bc_\alpha \in R$. 
Since $\sum bc_\alpha a_\alpha$ can be treated as an element of $M \subset V$, 
$$\sum c_\alpha a_\alpha = \frac{1}{b} \sum r_\gamma s_\gamma = \sum \frac{r_\gamma}{b} s_\gamma$$
for some $r_\gamma \in R$ with at most finitely many of them nonzero. 
Hence $S$ generates $V$ as a $K$-vector field. 

However, $\<S\> = V$ over $K$ does not entail $\<S\> = M$ over $R$. 
The case $M = \Z$ as a $\Z$-module, $V = \Q$ as a $\Q$-module, and $S = \{2\}$ serves as an easy counterexample .

\subsection{} 
Suppose $\Phi: F^{R}(A) \isom F^{R}(B)$ is an isomorphism. 
$A$ is a basis for $F^{R}(A)$, so it is particularly a maximal linearly independent set. 
Since the two free modules are isomorphic, $\Phi(A)$ is then a maximal linearly independent set for $F^{R}(B)$, which means $\#A = \# B$, i.e., $A \equiv B$ as sets. 

The reverse direction holds by universal property arguments. 


\end{document}
